\subsection{自由群的 p-滤子}

在本号中，\(p\) 表示一个素数，并假设 \(\mathbf{K} = \mathbf{F}_p\)。令 \(\gamma\) 为从 \(\mathrm{F}(\mathbf{X})\) 到 \(\Gamma(\mathbf{X})\) 的同态，它由对 \(\gamma(x) = 1 + x\) 中的 \(x\) 规定 \(\mathbf{X}\)；记 \(\mathrm{F}_n^{(p)}(\mathbf{X}) = \gamma(1 + \hat{\Lambda}_n(\mathbf{X}))\)。序列 \((\mathrm{F}_n^{(p)}(\mathbf{X}))_{n\geq 1}\) 是 \(\mathrm{F}(\mathbf{X})\) 上的一个整中心滤子，由于 \(\gamma\) 是单射（第 3 号，定理 1），它是分离的。它称为 \(p\) 上的 \(\mathrm{F}(\mathbf{X})\)-滤子。

命题 2。假设 X 是有限集。对每个整数 \(n \geqslant 1\)，群 \(F(X) / F_n^{(p)}(X)\) 是幂零类 \(p\) 的有限 \(\leqslant n\)-群。

按 \(n\) 归纳论证，只需证明对所有 \(\mathrm{F}_{n}^{(p)}(\mathbf{X}) / \mathrm{F}_{n + 1}^{(p)}(\mathbf{X})\)，\(p\) 是有限交换 \(n \geqslant 1\)-群。对所有 \(w \in \mathrm{F}_{n}^{(p)}(\mathbf{X})\)，元素 \(\gamma(w) - 1\) 属于 \(\hat{\mathbf{A}}(\mathbf{X})\)，且其阶 \(\geqslant n\)；记 \(\delta_n(w)\) 为 \(\gamma(w) - 1\) 的 n 次齐次分量。映射 \(n\) 是一个同态，其核为 \(\delta_n: \mathrm{F}_{n}^{(p)}(\mathbf{X}) \to \mathrm{A}^n (\mathbf{X})\)（§ 4，第 5 号，命题 3），因而 \(\mathrm{F}_{n + 1}^{(p)}(\mathbf{X})\) 同构于 \(\mathrm{F}_{n}^{(p)}(\mathbf{X}) / \mathrm{F}_{n + 1}^{(p)}(\mathbf{X})\) 的一个子群。由于 \(\mathrm{A}^n (\mathbf{X})\) 是有限集，\(\mathbf{X}\) 是 \(\mathrm{A}^n (\mathbf{X})\) 上的有限维向量空间，因而是有限交换 \(\mathbf{F}_p\)-群，所以 \(p\) 也是有限交换 \(\mathrm{F}_{n}^{(p)}(\mathbf{X}) / \mathrm{F}_{n + 1}^{(p)}(\mathbf{X})\)-群。

命题 3。对于 \(w \neq 1\) 中所有满足 \(\mathrm{F}(\mathbf{X})\) 的 w，存在一个有限 \(p\)-群 \(\mathbf{G}\) 以及一个从 \(f\) 到 \(\mathrm{F}(\mathbf{X})\) 的同态 \(\mathbf{G}\)，使得 \(f(w) \neq 1\)。

存在 \(x_{1},\ldots ,x_{r}\) 中的元素 \(\mathbf{X}\) 和整数 \(n_1,\dots ,n_r\)，使得 \(w = x_1^{n_1}\dots x_r^{n_r}\)。令 \(\mathrm{Y} = \{x_1,\dots,x_r\}\)。\(\mathrm{Y}\) 到 \(\mathbf{X}\) 的典范嵌入延拓为同态 \(\alpha : \mathrm{F}(\mathrm{Y})\to \mathrm{F}(\mathrm{X})\)；另一方面，令 \(\beta\) 为从 \(\mathrm{F}(\mathrm{X})\) 到 \(\mathrm{F}(\mathrm{Y})\) 的同态，其在 \(\mathbf{Y}\) 上的限制为恒等映射，并将 \(\mathbf{X} - \mathbf{Y}\) 映到 \(\{1\}\)。于是对 \(\beta (\alpha (y)) = y\)，有 \(y\in \mathbf{Y}\)，从而 \(\beta \circ \alpha\) 是 \(\mathrm{F}(\mathrm{Y})\) 的恒等自同构。显然存在 \(w^{\prime}\) 中的 \(\mathrm{F}(\mathrm{Y})\)，使得 \(w = \alpha (w^{\prime})\)；于是 \(\beta (w) = w^{\prime}\neq 1\)；现在 \(\bigcap_{n\geqslant 1}\mathbf{F}_n^{(p)}(\mathbf{Y}) = \{1\}\)，因此存在整数 \(n\geqslant 1\)，使得 \(\beta (w)\notin \mathbf{F}_{\mathrm{n}}^{(p)}(\mathbf{Y})\)。由命题 2，群 \(\mathbf{G} = \mathbf{F}(\mathbf{Y}) / \mathbf{F}_{\mathrm{n}}^{(p)}(\mathbf{Y})\) 是有限 \(p\)-群。若 \(f\) 是 \(\beta\) 与从 \(\mathrm{F}(\mathrm{Y})\) 到 \(\mathbf{G}\) 的典范同态的复合，则 \(f(w)\neq 1\)。

推论。F(X) 中所有有限指标正规子群的交为 \{1\}。

